\begin{derivation}[从\eqnrefs{eq:multimode_qpinem_continuous_kernel_dp101}到\eqnrefs{eq:multimode_qpinem_lattice_dp47}]

先把第 $\nu$ 个连续电子耦合算符记为
\[
\hat{\mathcal G}_{\nu,I}(t)
=
\sum_{ss'}
\int\dd^3p\,\dd^3p'\,
\mathcal G_{\nu;s's}(\mathbf p',\mathbf p;t)
|\mathbf p',s'\rangle\langle\mathbf p,s|.
\]
在实际窄带工作窗内，假定主导矩阵元集中在一个自旋保持、近固定动量转移通道
\[
(\mathbf p,s_0)
\longrightarrow
(\mathbf p+\hbar\boldsymbol\kappa_\nu,s_0).
\]
定义参考平移算符
\[
\hat b_\nu^\dagger
=
\int\dd^3p\,
|\mathbf p+\hbar\boldsymbol\kappa_\nu,s_0\rangle
\langle\mathbf p,s_0|
\]
以及投影后的共同速率 $\upsilon_\nu(t)$。于是精确耦合算符在该工作窗上分解成
\[
\hat{\mathcal G}_{\nu,I}(t)
=
\ii\upsilon_\nu(t)\hat b_\nu^\dagger
+
\hat R_{\nu,\kappa}(t).
\]
把所有余项收集为
\[
\hat R_\kappa(t)
=
\hbar
\sum_\nu
\left[
\hat R_{\nu,\kappa}(t)\hat a_\nu
+
\hat R_{\nu,\kappa}^\dagger(t)\hat a_\nu^\dagger
\right].
\]
令 $\hat H_{\rm tr}$ 只保留主导平移项，并由它生成 $\hat U_{\rm tr}$。完整传播子 $\hat U$ 与 $\hat U_{\rm tr}$ 的 Duhamel 恒等式为
\[
(\hat U-\hat U_{\rm tr})\Pi_{\rm in}
=
-\frac{\ii}{\hbar}
\int_{t_i}^{t_f}\dd t\,
\hat U(t_f,t)
\hat R_\kappa(t)
\hat U_{\rm tr}(t,t_i)\Pi_{\rm in}.
\]
完整传播子幺正，因此
\[
\|(\hat U-\hat U_{\rm tr})\Pi_{\rm in}\|
\le
\frac1\hbar
\int_{t_i}^{t_f}\dd t\,
\|
\hat R_\kappa(t)
\hat U_{\rm tr}(t,t_i)\Pi_{\rm in}
\|,
\]
这正是正文\eqnrefs{eq:multimode_momentum_step_projection_error_dp101}。只有该量足够小时，连续核才受控退化为
\[
\hat H_{\rm Q}^{(M)}
=
\ii\hbar\sum_\nu
\left[
\upsilon_\nu\hat b_\nu^\dagger\hat a_\nu
-\upsilon_\nu^*\hat b_\nu\hat a_\nu^\dagger
\right].
\]

设整数向量 $\boldsymbol\ell$ 在实际工作窗中唯一对应电子动量
\[
\mathbf p_{\boldsymbol\ell}
=
\mathbf p_0
+\hbar\sum_\nu\ell_\nu\boldsymbol\kappa_\nu,
\]
并定义
\[
|\boldsymbol\ell;\mathbf n\rangle
=
|\mathbf p_{\boldsymbol\ell}\rangle_e
\otimes
\bigotimes_\nu
|n_\nu-\ell_\nu\rangle_\nu .
\]
对第 $\nu$ 个吸收顶角，
\begin{align*}
\hat b_\nu^\dagger\hat a_\nu
|\boldsymbol\ell;\mathbf n\rangle
&=
\hat b_\nu^\dagger
|\mathbf p_{\boldsymbol\ell}\rangle
\otimes
\sqrt{n_\nu-\ell_\nu}\,
|n_\nu-\ell_\nu-1\rangle_\nu
\otimes_{\mu\neq\nu}
|n_\mu-\ell_\mu\rangle_\mu\\
&=
\sqrt{n_\nu-\ell_\nu}\,
|\boldsymbol\ell+\mathbf e_\nu;\mathbf n\rangle .
\end{align*}
同理，
\[
\hat b_\nu\hat a_\nu^\dagger
|\boldsymbol\ell;\mathbf n\rangle
=
\sqrt{n_\nu-\ell_\nu+1}\,
|\boldsymbol\ell-\mathbf e_\nu;\mathbf n\rangle .
\]

在旋去理想多频能量
$E_0+\sum_\nu\ell_\nu\hbar\omega_\nu$
的表象中，真实电子色散留下对角项
\[
\hat H_{\rm rec}
=
\hbar\sum_{\boldsymbol\ell}
\delta_{\boldsymbol\ell}
|\boldsymbol\ell\rangle\langle\boldsymbol\ell|.
\]
展开
\[
|\Psi_{\mathbf n}(t)\rangle
=
\sum_{\boldsymbol\ell}
c_{\boldsymbol\ell}^{(\mathbf n)}(t)
|\boldsymbol\ell;\mathbf n\rangle
\]
并代入
\[
\ii\hbar\partial_t|\Psi_{\mathbf n}\rangle
=
(\hat H_{\rm rec}+\hat H_{\rm Q}^{(M)})
|\Psi_{\mathbf n}\rangle.
\]
左乘 $\langle\boldsymbol\ell;\mathbf n|$ 后，来自相邻格点的两类项分别为
\[
+\ii\hbar
\upsilon_\nu
\sqrt{n_\nu-\ell_\nu+1}\,
c_{\boldsymbol\ell-\mathbf e_\nu}^{(\mathbf n)}
\]
和
\[
-\ii\hbar
\upsilon_\nu^*
\sqrt{n_\nu-\ell_\nu}\,
c_{\boldsymbol\ell+\mathbf e_\nu}^{(\mathbf n)}.
\]
与对角反冲项合并并除以 $\hbar$，得到正文\eqnrefs{eq:multimode_qpinem_lattice_dp47}。

最后说明多指标标签为什么需要一一对应。若存在
$\boldsymbol\ell\neq\boldsymbol\ell'$ 却有
\[
\sum_\nu
(\ell_\nu-\ell_\nu')
\boldsymbol\kappa_\nu=0,
\]
则
$|\mathbf p_{\boldsymbol\ell}\rangle=
|\mathbf p_{\boldsymbol\ell'}\rangle$，
这两个“格点”并非正交电子态。此时必须先把所有到达同一物理动量的路径振幅相加，再定义独立基底；直接把 $\boldsymbol\ell,\boldsymbol\ell'$ 当成两个末态会重复计算 Hilbert 空间维数。
\end{derivation}

\begin{derivation}[从\eqnrefs{eq:multicolor_classical_scattering_dp47}到\eqnrefs{eq:multicolor_pinem_convolution_dp47}]

在多模相干态经典极限中，每个量子模式的算符被分解成平均幅与量子涨落；固定
\[
\beta_\nu=\alpha_\nu g_{\\mathrm Q,\nu}
\]
并令相对涨落修正消失后，电子散射算符为
\[
S_{\rm cl}^{(M)}
=
\exp\!\left[
\sum_\nu
(\beta_\nu\hat b_\nu^\dagger-\beta_\nu^*\hat b_\nu)
\right].
\]
完整连续动量空间中的平移算符彼此对易，
\[
[\hat b_\mu,\hat b_\nu]=
[\hat b_\mu,\hat b_\nu^\dagger]=0,
\]
所以指数可逐模分解：
\[
S_{\rm cl}^{(M)}
=
\prod_\nu
\exp(
\beta_\nu\hat b_\nu^\dagger
-\beta_\nu^*\hat b_\nu).
\]

对单个模式，记
\[
\beta_\nu=|\beta_\nu|\ee^{\ii\varphi_\nu},
\qquad
\hat T_\nu\equiv\hat b_\nu^\dagger.
\]
在理想双边平移空间中 $\hat T_\nu$ 是幺正算符，
\[
\hat T_\nu^{-1}=\hat b_\nu.
\]
因此任意整数幂统一写作
\[
\hat T_\nu^{r}
=
\begin{cases}
(\hat b_\nu^\dagger)^r,&r\ge0,\\
(\hat b_\nu)^{-r},&r<0.
\end{cases}
\]
取 $\hat T_\nu$ 的相位本征态
\[
\hat T_\nu|\theta_\nu\rangle
=
\ee^{-\ii\theta_\nu}|\theta_\nu\rangle,
\]
则单模生成元在该表象中成为普通复数：
\begin{align*}
\beta_\nu\hat T_\nu-\beta_\nu^*\hat T_\nu^{-1}
&\longrightarrow
|\beta_\nu|
\left[
\ee^{\ii(\varphi_\nu-\theta_\nu)}
-\ee^{-\ii(\varphi_\nu-\theta_\nu)}
\right]\\
&=
2\ii|\beta_\nu|
\sin(\varphi_\nu-\theta_\nu).
\end{align*}
Jacobi--Anger 恒等式给
\[
\exp[
2\ii|\beta_\nu|\sin(\varphi_\nu-\theta_\nu)
]
=
\sum_{r_\nu\in\mathbb Z}
J_{r_\nu}(2|\beta_\nu|)
\ee^{\ii r_\nu\varphi_\nu}
\ee^{-\ii r_\nu\theta_\nu}.
\]
把
$\ee^{-\ii r_\nu\theta_\nu}$
重新提升为算符 $\hat T_\nu^{r_\nu}$，得到严格的整数平移展开
\[
\exp(
\beta_\nu\hat b_\nu^\dagger
-\beta_\nu^*\hat b_\nu)
=
\sum_{r_\nu\in\mathbb Z}
J_{r_\nu}(2|\beta_\nu|)
\ee^{\ii r_\nu\varphi_\nu}
\hat T_\nu^{r_\nu}.
\]

对所有模式相乘：
\[
S_{\rm cl}^{(M)}
=
\sum_{\{r_\nu\}}
\left[
\prod_\nu
J_{r_\nu}(2|\beta_\nu|)
\ee^{\ii r_\nu\varphi_\nu}
\right]
\prod_\nu\hat T_\nu^{r_\nu}.
\]
若
\[
\kappa_\nu=m_\nu\kappa_0,
\qquad
m_\nu\in\mathbb Z,
\]
则
\[
\prod_\nu\hat T_\nu^{r_\nu}
\]
把电子总共平移
\[
\Delta p
=
\hbar\kappa_0
\sum_\nu m_\nu r_\nu.
\]
因此所有满足
\[
\ell=\sum_\nu m_\nu r_\nu
\]
的整数路径都到达同一个电子末态 $|\ell\rangle$，其振幅必须相干求和。于是
\[
A_\ell
=
\sum_{\{r_\nu\}}
\delta_{\ell,\sum_\nu m_\nu r_\nu}
\prod_\nu
J_{r_\nu}(2|\beta_\nu|)
\ee^{\ii r_\nu\varphi_\nu},
\]
即正文\eqnrefs{eq:multicolor_pinem_convolution_dp47}。

若不同颜色的动量步长不满足这种整数公度关系，则一般不存在唯一的一维总阶数 $\ell$；此时末态应继续用真实电子动量或多指标通道标记，不能把不同可分辨末态强行卷积到一条 Bessel 能量格。
\end{derivation}

